\section{Schema \& Data Model}
InfluxDB organises data as:
\begin{center}
Organisation --> bucket --> Measurement --> Tags + Field + Timestamp
\end{center}

\subsection{Organisation}
A single InfluxDB organisation \texttt{qut-microgrid} is used, consistent with the Telegraf output configuration in the MQTT design document.

\subsection{Buckets}

\begin{table}[H]
  \caption{Buckets}
  \begin{center}
    \begin{tabular}{lll}
      \hline
      \textbf{Bucket} & \textbf{Purpose} & \textbf{Retention} \\
      \hline
      \texttt{telemetry} & Full--resolution samples from all nodes & 90 days (TBD) \\
      \texttt{telemetry\_1m} & 1--minute aggregates for long-term storage & Indefinite \\
      \hline
    \end{tabular}
  \end{center}
\end{table}

The bucket name \texttt{telemetry} matches the Telegraf output configuration. A seperate downsampled bucket avoids data retention growth while preserving long--term trend visibility. 
      
\subsection{Measurements}
Each MQTT subsystem segment maps to a distinct InfluxDB measurement name:

\begin{table}[H]
  \caption{Measuremnts}
  \begin{center}
    \begin{tabular}{ll}
      \hline
      \textbf{Measurement} & \textbf{Description}\\
      \hline
      \texttt{voltage} & Phase voltage samples (V) \\
      \texttt{current} & Phase and neutral current samples (A) \\
      \texttt{status} & Node liveness and heartbeat dignostics \\
      \texttt{ptp\_status} & PTP synchronisation diagnostics \\
      \hline
    \end{tabular}
  \end{center}
\end{table}

Additional subsystems may be introduced without structural schema changes if the STM32 firmware is extended to publish derived quantities.

\subsection{Tags}
Tags are string--valued, indexed, and low--cardinality. All three tags are populated automatically by Telegraf's \texttt{topic\_parsing} from the four--level MQTT topic hierachy. 

\begin{table}[H]
  \caption{Tags}
  \begin{center}
    \begin{tabular}{llp{0.2\linewidth}p{0.4\linewidth}}
      \hline
      \textbf{Tag Key} & \textbf{Source} & \textbf{Example Values} & \textbf{Notes}\\
      \hline
      \texttt{node\_id} & Topic segment 2 & \texttt{node\_1}, \texttt{node\_2} & Fixed at firmware flash time, MAC/IP addresses explicitly avoided. \\
      \texttt{measurement} & Topic segment 4 & \texttt{va}, \texttt{ib} & Aligned with IEEE-1459-2010 notation \\
      \hline
    \end{tabular}
  \end{center}
\end{table}

\subsection{Fields}
\begin{table}[H]
  \caption{Tags}
  \begin{center}
    \begin{tabular}{lllp{0.5\linewidth}}
      \hline
      \textbf{Field Key} & \textbf{Type} & \textbf{Source} & \textbf{Units}\\
      \hline
      \texttt{v} & float64 & JSON payload \texttt{v} field & implied by \texttt{\_measurement} and \texttt{measurement} tag (V or A) \\
      \hline
    \end{tabular}
  \end{center}
\end{table}

\subsection{Timestamp}
Timestamps are written at nanosecond precision using the \texttt{ts} field from the JSON payload, sourced from the STM32 IEEE 1588 hardware timestamping peripheral. The historian does not substitute broker receive time or Telegraf processing time: the node--reported PTP timestamp is authoritative end--to--end.

\subsection{Representative Schema}
\begin{table}[H]
  \caption{Representative Schema}
  \begin{center}
    \begin{tabular}{lllll}
      \hline
      \textbf{subsystem (\_measurement)} & \textbf{node\_id} & \textbf{measurement} & \textbf{\_field} & \textbf{Example Value}\\
      \hline
      \texttt{voltage} & \texttt{node\_1} & \texttt{va} & \texttt{v} & \\
      \texttt{voltage} & \texttt{node\_1} & \texttt{vb} & \texttt{v} & \\
      \texttt{voltage} & \texttt{node\_1} & \texttt{vc} & \texttt{v} & \\
      \texttt{current} & \texttt{node\_1} & \texttt{ia} & \texttt{v} & \\
      \texttt{current} & \texttt{node\_1} & \texttt{ib} & \texttt{v} & \\
      \texttt{current} & \texttt{node\_1} & \texttt{ic} & \texttt{v} & \\
      \texttt{ptp\_status} & \texttt{node\_1} & \texttt{offset\_ns} & \texttt{v} & \\
      \hline
    \end{tabular}
  \end{center}
\end{table}







